%% Artificial Life (MIT Press) journal submission
%% MUST COMPILE WITH XELATEX OR LUALATEX

\documentclass{artificial-life}

\usepackage[style=apa,natbib=true]{biblatex}
\addbibresource{Le.bib}

\title{A collective-level Darwinian population
without collective-level selection}

\auth{
    Nam H. Le \affil{1}
}

\corresponding{Nam H. Le}{namlehai90@gmail.com}

\affiliations{
    \item Independent researcher
}

\abs{
Evolutionary transitions in individuality gave life its layered structure: free
cells became the parts of organisms, solitary insects the parts of colonies. How
a new level of individual first arises is difficult to model, because the
standard accounts of group and kin selection begin by assuming a group and a
criterion that evaluates it. Artificial life can avoid that circularity by
writing no fitness function at all, but the systems built this way have
generally produced a replicating pattern rather than a population of
higher-level entities standing in relations of descent. Here I report such a population in VitaNova, a
minimal model derived from flocking rules and an energetic ecology, in which it
appears without being specified. Collectives form, grow and divide,
and the divisions compound: across sixteen independent runs, several hundred
collectives arose, a third of them divided, and the deepest lineage ran nineteen
collective generations. Experimental results show that this population meets all
three requirements for a Darwinian population at the collective level.
Collectives reproduce; daughters resemble their parents in size across the
lineage, with a parent--offspring regression slope near two thirds that survives
controls for replicate and for temporal phase; and reproductive success varies
sharply, with three fifths of collectives leaving no offspring at all. Nothing
in the model refers to a collective, so none of this was imposed. Ablations
locate what the population depends upon: an ecology in which grouping pays, and
unbounded growth, since a size ceiling that stabilises the collectives'
internal structure collapses the deepest lineage to a single generation. It does
not depend on evolution, on relatedness, on a division of labour, or on fitness
decoupling. What is absent is selection acting upon the collectives themselves,
because the collective property that is heritable and responds to selection is
internal structure, while the rate at which collectives divide does not respond
at all. The system reaches the starting line of a transition, which theory holds
to be the hard part, and stops there. The measurements show where the next
barrier lies.
}

\keywords{evolutionary transitions in individuality, collective reproduction,
Darwinian populations, self-organisation, collective motion, agent-based model}

\begin{document}

% Don't Touch ! %%%%%%
\coverpage           %
\linenumbers         %
\doublespacing       %
%%%%%%%%%%%%%%%%%%%%%%

\section{Introduction}

Evolution does not only adapt individuals to their environments. It also makes
new kinds of individual out of old ones. Free-living cells became the parts of
multicellular organisms, solitary insects became the parts of colonies, and
separate replicators became the parts of chromosomes. These events, the
evolutionary transitions in individuality
\citep{smith1995major,west2015major}, are the reason biological complexity is
layered rather than flat.

Modelling them runs into a structural difficulty. Selection acts on entities
that already reproduce differentially, so any account of how a new reproducing
entity comes into being is liable to assume what it should explain. The standard
formal treatments begin by specifying a fitness at the collective level, or a
relatedness among a group's members, and then ask when cooperation is favoured
\citep{okasha2005multilevel,hamilton1964genetical}. Those are the right tools
for asking how an existing group behaves, and the wrong ones for asking where
the group came from.

Artificial life has a standard way around the circularity, which is to write no
fitness function at all and let differential reproduction fall out of an ecology
\citep{ray1991approach,yaeger1994polyworld,channon2001geb}. What the field has
mostly produced under that stance, however, is \emph{pattern replication}:
configurations that reproduce their own form
\citep{langton1984self,sayama1999evoloop,chan2020lenia}. A replicating pattern
is not a Darwinian population. For that one needs many higher-level entities,
standing in relations of descent, differing in what they transmit and in how
many descendants they leave \citep{godfreysmith2009}.

The system studied here, VitaNova, was introduced in an earlier paper
\citep{le2025vitanova}, which described its collectives in simple terms: they
form, acquire a ring-like internal structure, and sometimes divide. Nothing
about collective-level descent or heredity was measured there. This article
supplies those measurements.

Here I report that the population is real. Across sixteen independent runs the
divisions compound into lineages up to nineteen collective generations deep, a
collective-level trait is transmitted across those lineages, and reproductive
success varies sharply among collectives. Nothing in the model refers to a
collective, and no collective-level selection is imposed or possible.

The article then asks what the population depends upon, by switching components
off, and what it still lacks. Together these locate the system: at the starting
line of a transition, which theory holds to be the hard part, with the next
barrier visible.

\section{Background}
\label{sec:related}

\subsection{What a transition is held to require}

\citet{godfreysmith2009} sets out three requirements for a Darwinian
population: its members reproduce, offspring resemble their parents, and that
resemblance covaries with reproductive success. Applied one level up, these
become three questions. Do collectives reproduce? Is a collective-level property
transmitted from a collective to its offspring collectives? Do collectives
differ in how many offspring they leave? The requirements are modest, and that
is their value: they are stated so as to be checkable.

Around them sits a body of theory concerned with mechanisms that are held to
enable a transition or to mark one. \citet{roze2001mutation} treat a single-cell
bottleneck as the device that makes collective-level heredity possible, by
purging the within-collective variation that would otherwise decouple a
collective from its offspring. \citet{michod2003reorganization} propose fitness
decoupling, the collective's fitness ceasing to be the average of its parts', as
the signature that a transition has occurred; \citet{bourrat2022tradeoff} reply
that apparent decoupling is an artefact of a shifting reference environment and
propose tradeoff breaking in its place. \citet{queller1997cooperators}
distinguishes fraternal transitions among relatives, where relatedness aligns
interests, from egalitarian transitions among unlike partners, and
\citet{tarnita2013evolution} draw the related distinction between collectives
that stay together after dividing and collectives that come together from
dispersed parts. \citet{grosberg2007evolution} review how these considerations
play out in the origin of multicellularity, and \citet{boomsma2009lifetime}
argues that lifetime monogamy, and hence high relatedness, is what made
eusociality possible.

Two features of this literature bear on what follows. The first is that its
central quantities are usually assessed qualitatively: a system is described,
the reader is invited to agree that its collectives reproduce, and the criteria
are satisfied by narration rather than by measurement. The second is that
almost all of it concerns what happens to a collective-level population once one
exists. The question of how such a population arises in the first place, before
anything can select upon it, is precisely the question the circularity makes
hard to pose.

\subsection{The circularity, and implicit fitness}

The circularity is avoided in artificial life by refusing to state an objective,
and the practice is long established. In Tierra \citep{ray1991approach},
self-replicating programs compete for processor time and for memory, and their
differential replication is a consequence of what they do rather than of a score
assigned to them. In Polyworld \citep{yaeger1994polyworld}, agents carrying
evolvable neural controllers must find food and mate in order to leave
descendants, and nothing in the simulator evaluates how well they are doing.
\citet{channon2001geb} makes the argument explicit, that open-ended evolution
requires exactly this refusal, and builds Geb upon it.

The stance has two consequences, and both matter here. Because no criterion is
written down, no property of the system can have been produced by one, so a
collective that appears is not an artefact of a specification. That is the
whole point. The cost is that nothing can be steered: one cannot select for a
collective-level property, or arrange a bottleneck, or impose a life cycle. It
is for this reason that implicit-fitness systems are seldom put forward as
models of transitions, and for the same reason that a collective-level Darwinian
population, should one appear in such a system, is worth measuring carefully.

\subsection{What artificial and experimental systems have shown}

Table~\ref{tab:systems} sets the main systems, VitaNova among them, against
\citet{godfreysmith2009}'s three requirements, and the pattern in the table is
that collective reproduction and collective heredity seldom arise together
unless one of them has been supplied.

Self-reproducing structures in cellular automata, from \citet{langton1984self}
through the evoloop of \citet{sayama1999evoloop} to the continuous patterns of
\citet{chan2020lenia}, replicate a configuration beautifully and without any
population of heritable variants competing on fitness. The components have no
heredity of their own and nothing selects among them, so there is no lower level
from which a higher level might arise. The digital evolution platforms have the
opposite profile: Tierra \citep{ray1991approach} and Avida
\citep{ofria2004avida} have genuine heredity and genuine selection, at one level
only, with no collective that reproduces as a unit. DISHTINY
\citep{moreno2019dishtiny} was designed to host hierarchical individuality and
does so, by providing explicit group-level machinery, which answers a question
about sufficiency rather than about origin.

Experimental evolution has made the most concrete progress on transitions
themselves. The snowflake yeast of \citet{ratcliff2012experimental} evolve
simple multicellularity within weeks under settling selection, and
\citet{hammerschmidt2014life} evolve a life cycle in bacterial mats.
\citet{black2020scaffolding} state the underlying logic most clearly: an
externally imposed ecological scaffold is \emph{sufficient} to confer Darwinian
properties on collectives that do not possess them intrinsically. The common
feature of these designs, and their limit for the present question, is that
selection at the collective level is supplied by the experimenter. They
establish what collective-level selection can build. They do not establish how a
collective-level Darwinian population comes to exist before anything selects
upon it.

\begin{table}[tp]
\centering
\small
\begin{tabular}{llll}
\hline
System & Collective & Collective & Selection at the \\
       & reproduces & heredity & collective level \\
\hline
Self-reproducing CA \citep{langton1984self} & yes & no & no \\
Evoloop \citep{sayama1999evoloop} & yes & partial & implicit \\
Lenia \citep{chan2020lenia} & yes & no & no \\
Tierra, Avida \citep{ray1991approach,ofria2004avida} & no & n/a & no \\
Polyworld \citep{yaeger1994polyworld} & no & n/a & no \\
DISHTINY \citep{moreno2019dishtiny} & yes & yes & designed in \\
Scaffolded collectives \citep{black2020scaffolding} & imposed & yes & imposed \\
Snowflake yeast \citep{ratcliff2012experimental} & yes & yes & imposed \\
VitaNova (this article) & \textbf{yes} & \textbf{yes} & \textbf{none, and none possible} \\
\hline
\end{tabular}
\caption{Systems cited in discussions of artificial transitions, set against the
three requirements of \citet{godfreysmith2009}. ``Imposed'' marks a property
supplied by the experimental protocol rather than arising from the system's own
dynamics. The final row is the present article; the paper that introduced
VitaNova \citep{le2025vitanova} reported neither collective heredity nor
collective descent.}
\label{tab:systems}
\end{table}

\subsection{What is not measured}

Two gaps follow from the above, and they compound each other.

The first is that collective-level descent is rarely followed. A system is
reported to exhibit group reproduction on the strength of division events, but a
division event and a lineage are different things: a hundred divisions that
never compound describe a very different system from a hundred that reach
nineteen generations of descent, and only the second is a population on which
selection could act. Reporting the depth of collective lineages, rather than the
count of divisions, is cheap and rarely done.

The second is that collective-level heredity is rarely measured against the
controls it needs. A correlation between a parent collective and its offspring
is the expected consequence of a shared temporal trend, or of some runs simply
producing larger collectives than others, and separating transmission from
either requires centring within replicate and within phase and testing against a
permutation null that preserves that structure. The measurement tradition in
this field, represented by \citet{dolson2019modes}, tracks novelty and
complexity rather than descent and transmission, so the apparatus is not to
hand.

This article supplies both for VitaNova: the depth of its collective lineages,
and the transmission of a collective-level trait across them under a ladder of
fixed-effects controls. Neither quantity was reported when the system was
introduced.


\section{VitaNova}
\label{sec:model}

\subsection{Agents and their controllers}

The world is an $800 \times 600$ px arena with reflecting boundaries, holding
$100$ agents and $10$ predators at initialisation (Figure~\ref{fig:system}a).
Each agent perceives others within a radius $R = 25$ px, and the separation,
alignment and cohesion radii of the underlying flocking rules
\citep{reynolds1987flocks} are set to the same value
(Figure~\ref{fig:system}b). That the four radii coincide is a property of the
original implementation rather than a design choice, and
Section~\ref{sec:conditions} shows it has consequences.

Each agent carries a genome of $170$ real-valued weights specifying a
feed-forward network with eleven inputs, ten hidden units and six outputs. The
inputs are the relative positions of three perceived neighbours, the relative
position of the nearest perceived predator, the agent's own velocity, and a
local density estimate. The outputs are not motions but \emph{tendencies}: four
gains in $[-1,1]$, produced by a $\tanh$, that weight the separation, alignment,
cohesion and predator-avoidance terms, together with two further outputs that
set a worker or soldier state. Writing $\tau_s, \tau_a, \tau_c, \tau_p$ for the
four gains, $\mathbf{s}, \mathbf{a}, \mathbf{c}$ for the three flocking vectors
and $\mathbf{q}$ for the displacement to the nearest perceived predator, the
velocity update is
\begin{equation}\label{eqn:update}
\mathbf{v} \;\leftarrow\; \sigma \cdot \widehat{\left(
\mathbf{v} + \tau_s \mathbf{s} + \tau_a \mathbf{a} + \tau_c \mathbf{c}
- \max(\tau_p, 0)\,\frac{\mathbf{q}}{\|\mathbf{q}\|^{2}}\right)},
\end{equation}
with $\sigma$ a fixed speed and the hat denoting normalisation. The three
flocking gains are signed, so an agent may invert any of them; the avoidance
gain is rectified, so an agent cannot be attracted to a predator. An agent's
controller therefore chooses a style of flocking rather than a trajectory, and
what a collective looks like is a consequence of the styles its members happen
to have.

Two details are inherited from the original implementation and preserved rather
than corrected: alignment and cohesion are computed including the focal agent
while separation excludes it, and the three neighbours entering the input vector
are taken in array index order rather than by distance.

\subsection{Energy, reproduction, and the absence of a fitness function}
\label{sec:energy}

An agent gains $1.1$ energy per step and pays $1.0$. On top of this it receives
$+0.2$ when between two and nineteen other agents lie within $15$ px, and $-0.2$
at twenty or more. Reproduction is asexual and occurs when energy exceeds $200$;
it costs $100$, places the offspring nearby, and mutates its genome with rate
and standard deviation $0.1$. Initial energy is $50$. Death is by age alone, at
$1000$ steps, so energy gates reproduction but never kills.

Following the implicit-fitness stance of \citet{ray1991approach} and
\citet{yaeger1994polyworld}, there is no fitness function anywhere in VitaNova,
no term in any agent's accounting that refers to a collective, and no rule that
evaluates one. Differential reproduction arises because agents in company gain
energy faster than agents alone, which is a property of the ecology rather than
a criterion imposed upon it. This is what makes the collectives reported below
worth measuring: since no group criterion is written down, they cannot be an
artefact of one.

The constants above have a consequence that is easier to see by multiplying them
out than by running the model, and Section~\ref{sec:conditions} turns on it. An
agent that never receives the grouping subsidy accumulates at most
$50 + 0.1 \times 1000 = 150$ energy in an entire lifetime, against a
reproduction threshold of $200$. It cannot reproduce under any behaviour, for
any seed. Grouping in VitaNova is therefore not an advantage that agents
discover; it is an arithmetic precondition for reproducing at all.

Agents also switch between worker and soldier states according to the two
remaining network outputs, with a small difference in energy accounting and a
predator rule that responds to groups of soldiers. Ablating this caste machinery
makes no difference to anything reported here, and I leave its details to a
separate report.

\subsection{Implementation and validation}

All experiments use a vectorised reimplementation of the original, which is what
makes the ablations affordable: it runs $90$ to $110$ times faster. It is not
trajectory-exact, because the original updates agents in sequence, so that an
agent sees neighbours that have already moved within the step, while the port
updates every agent from a single snapshot. Validation is therefore
distributional, and on that basis the two agree. Over $78$ seeds the population
growth-rate ratio is $1.04$, a collective divided in $32\%$ of runs against
$30\%$, and the size of the dividing collective does not differ
(Mann--Whitney $p = 0.35$).

\begin{figure}[tp]
    \centering
    \includegraphics[width=\textwidth]{Le.fig01.pdf}
    \caption{VitaNova, and one run of it. (a) The $800 \times 600$ px arena at
    step 900 of one run, holding 73 agents and 10 predators. Agents are coloured
    by state, workers in blue and soldiers in amber; predators are dark red
    triangles; short grey arrows are velocities. This colour scheme is used for
    agents throughout. (b) The world of one agent, centred on the agent with the
    most neighbours, which is ringed. The outer dashed circle is the radius
    $R = 25$ px at which perception, separation, alignment and cohesion all act;
    the inner dotted circle is the $15$ px range over which the energy subsidy
    counts neighbours. (c--f) The largest collective at four times, each drawn
    in its own frame with the collective on a shaded disc; in (f) the two
    products of a division are shaded separately. (g) Annularity, the measure of
    internal structure defined in Section~\ref{sec:measure} (left axis, solid),
    and the size of the largest collective (right axis, dashed, logarithmic),
    over the whole run.}
    \label{fig:system}
\end{figure}

\end{document}
